\Askhsh{13751}{4}
\begin{ylh}{1}
    \item 3.6. Κριτήρια ισότητας ορθογώνιων τριγώνων
    \item 5.6. Εφαρμογές στα τρίγωνα.
\end{ylh}
% ===================================================================
%  Αρχείο: 13751.tex (Fragment)
%  Αυτόματη μετατροπή μέσω doc2latex.py
% ===================================================================

ΘΕΜΑ 4

Στο παρακάτω σχήμα οι ευθείες ε\textsubscript{1} και ε\textsubscript{2} είναι παράλληλες. Το τρίγωνο ΑΒΓ είναι οξυγώνιο, ενώ το ΔΕΖ είναι αμβλυγώνιο με $\widehat{Ε}$ \textgreater{} 90\textsuperscript{ο} . Ισχύει επίσης ότι ΑΓ = ΔΖ.

\begin{alist}
\item

\begin{alist}
\item
  Να σχεδιάσετε τα ύψη των τριγώνων από τις κορυφές Α και Δ ονομάζοντάς τα ΑΗ και ΔΘ αντίστοιχα. \monades{05}
\item
  Να αποδείξετε ότι ΗΓ = ΘΖ. \monades{12}

\item Να δικαιολογήσετε γιατί ΕΖ \textless{} ΒΓ. \monades{08}


\begin{center}
\begin{tikzpicture}[scale=1]
\tkzDefPoint(-1,3){L1_left}
\tkzDefPoint(8,3){L1_right}
\tkzDefPoint(-1,0){L2_left}
\tkzDefPoint(8,0){L2_right}

\tkzDefPoint(1,3){A}
\tkzDefPoint(0,0){B}
\tkzDefPoint(2.5,0){C}
\tkzDefPoint(5,3){D}
\tkzDefPoint(4.5,0){E}
\tkzDefPoint(6.5,0){Z}

\tkzDefPointBy[projection=onto L2_left--L2_right](A) \tkzGetPoint{H}
\tkzDefPointBy[projection=onto L2_left--L2_right](D) \tkzGetPoint{Theta}

\draw[l3] (A)--(B)--(C)--cycle;
\draw[l3] (D)--(E)--(Z)--cycle;
\draw[l3, dashed] (A)--(H);
\draw[l3, dashed] (D)--(Theta);

\draw[l3] (L1_left)--(L1_right) node[right] {$\varepsilon_1$};
\draw[l3] (L2_left)--(L2_right) node[right] {$\varepsilon_2$};

\tkzDrawPoints(A,B,C,D,E,Z,H,Theta)
\tkzLabelPoint[above](A){$A$}
\tkzLabelPoint[below](B){$B$}
\tkzLabelPoint[below](C){$\Gamma$}
\tkzLabelPoint[above](D){$\Delta$}
\tkzLabelPoint[below](E){$E$}
\tkzLabelPoint[below](Z){$Z$}
\tkzLabelPoint[below](H){$H$}
\tkzLabelPoint[below](Theta){$\Theta$}

\tkzMarkRightAngle(A,H,C)
\tkzMarkRightAngle(D,Theta,E)
\end{tikzpicture}
\end{center}
\end{alist}
\end{alist}
